Lors d'une expérience de diffusion Compton, des rayons X de \qty{\v{l0}}{pm} frappent une cible de graphite. Un détecteur placé à un certain angle mesure des photons diffusés de \qty[round-mode=places,round-precision=3]{\fpeval{\v{l0} + 1e12*6.626e-34/(9.109e-31*3e8)*(1 - cosd(\v{th}))}}{pm}.

\begin{enumerate}[label=\alph*)]
    \item À quel angle le détecteur est-il placé ?
    \item Quelle énergie cinétique (en eV) l'électron a-t-il reçue ?
    \item Le détecteur mesure aussi des photons de \qty{\v{l0}}{pm}, non décalés. D'où viennent-ils ?
\end{enumerate}
